\documentclass[11pt]{article}
\usepackage{mathblatt}
% Baustein im Vorspann (fehlt in der Anleitung): Rechteck mit Rand, für Sachaufgaben „Beet mit Weg"
% #1 Breite innen, #2 Länge innen, #3 Randbreite (cm), #4 Label Breite, #5 Label Länge, #6 Text im Rand
\newcommand{\rechteckrand}[6]{%
\begin{tikzpicture}
\fill[gray!18] (0,0) rectangle ({#2+2*#3},{#1+2*#3});
\fill[white] (#3,#3) rectangle ({#2+#3},{#1+#3});
\draw[thick] (0,0) rectangle ({#2+2*#3},{#1+2*#3});
\draw[thick] (#3,#3) rectangle ({#2+#3},{#1+#3});
\node[above] at ({#3+#2/2},#3) {$#5$};
\node[right] at (#3,{#3+#1/2}) {$#4$};
\node at ({#3+#2/2},{#1+1.5*#3}) {\small #6};
\end{tikzpicture}}
\begin{document}
\blattkopf{Quadratische Gleichungen}{Lernblatt}
\weit
\verzeichniszeile{\verz{e1}{1 Wurzelziehen und Lösbarkeit (Nr. 12–23)} \verztrenn \verz{e2}{2 Satz vom Nullprodukt (Nr. 24–31)} \verztrenn \verz{e3}{3 Normalform und p-q-Formel (Nr. 32–44)} \verztrenn \verz{e4}{4 Sachaufgaben (Nr. 45–49)} \verztrenn \verz{abhaken}{Das kann ich}}
% Einheit 1 – Wurzelziehen und Lösbarkeit – Nr. 12–23
\setcounter{aufgabe}{11}
\einheitenkopf[e1]{Einheit 1 von 4 · Wurzelziehen und Lösbarkeit}
\zweigzeile{Hier lernst du, Gleichungen mit $x^2$, aber ohne $x$-Glied durch Wurzelziehen zu lösen und zu sagen, wie viele Lösungen sie haben · neu oder schon bekannt – je nach Buch (an der Oberschule je nach Buch Kl. 9 oder 10) · P10 · baut auf: Quadrieren, Wurzelziehen, lineare Gleichungen, Scheitelpunktform}

\begin{aufgabe}{Ich erkenne, ob eine Gleichung quadratisch ist.}
Rechne nicht. Kreuze an, ob die Gleichung quadratisch oder linear ist.
\begin{teile}
\teil $3x + 4 = 10$ \hfill \kreuz{quadratisch} \quad \kreuz{linear}
\teil $x^2 = 50$ \hfill \kreuz{quadratisch} \quad \kreuz{linear}
\teil $(x - 1)^2 = 9$ \hfill \kreuz{quadratisch} \quad \kreuz{linear}
\teil $2\cdot(x + 3) = 8$ \hfill \kreuz{quadratisch} \quad \kreuz{linear}
\teil $5x = x^2$ \hfill \kreuz{quadratisch} \quad \kreuz{linear}
\end{teile}
\end{aufgabe}

\begin{aufgabe}{Ich erkenne an der Zahl rechts, wie viele Lösungen eine Gleichung wie $x^2 = 5$ hat.}
Rechne nicht. Kreuze an, wie viele Lösungen die Gleichung hat.
\begin{teile}
\teil $x^2 = 16$ \hfill \kreuz{zwei} \quad \kreuz{eine} \quad \kreuz{keine}
\teil $x^2 = -9$ \hfill \kreuz{zwei} \quad \kreuz{eine} \quad \kreuz{keine}
\teil $x^2 = 2$ \hfill \kreuz{zwei} \quad \kreuz{eine} \quad \kreuz{keine}
\teil $5x^2 = 0$ \hfill \kreuz{zwei} \quad \kreuz{eine} \quad \kreuz{keine}
\teil $x^2 = -0{,}5$ \hfill \kreuz{zwei} \quad \kreuz{eine} \quad \kreuz{keine}
\end{teile}
\end{aufgabe}

\verfahren{Gleichungen der Form $x^2 = $ Zahl}

\begin{aufgabe}{Ich kann eine Gleichung der Form $x^2 = $ Zahl durch Wurzelziehen lösen.}
Löse die Gleichung.
\rechenplatz{2}
\begin{gleichungsraster}[2]
\gl{x^2 = 25} & \gl{x^2 = 81} \\
\gl{x^2 = 4} & \gl{x^2 = 144} \\
\anweisung{Rechne mit dem Taschenrechner, wenn die Wurzel nicht aufgeht. Gib die Lösungen mit Wurzel und auf zwei Stellen nach dem Komma gerundet an.}
\gl{x^2 = 11} & \gl{x^2 = -16} \\
\end{gleichungsraster}
\end{aufgabe}

\begin{aufgabe}{Ich kann erst $x^2$ freistellen und dann die Wurzel ziehen.}
Löse die Gleichung.
\rechenplatz{3}
\begin{gleichungsraster}[3]
\gl{x^2 + 3 = 19} & \gl{x^2 - 20 = 5} \\
\gl{5x^2 = 45} & \gl{3x^2 + 5 = 32} \\
\gl{4x^2 - 9 = -9} & \gl{2x^2 + 11 = 3} \\
\gl{3x^2 - 4 = 17} & \\
\end{gleichungsraster}
\end{aufgabe}

\verfahren{Gleichungen mit quadrierter Klammer}

\begin{aufgabe}{Ich kann eine Gleichung mit quadrierter Klammer durch Rückwärtsrechnen lösen.}
Löse die Gleichung.
\rechenplatz{3}
\begin{gleichungsraster}[3]
\gl{(x - 2)^2 = 9} & \gl{(x - 1)^2 = 16} \\
\gl{(x - 5)^2 = 4} & \gl{(x - 3)^2 = 25} \\
\gl{(x + 3)^2 = 49} & \gl{(x + 6)^2 = 0} \\
\gl{(x - 2)^2 = 5} & \\
\end{gleichungsraster}
\end{aufgabe}

\verfahren{Lösungen prüfen und zählen}

\begin{aufgabe}{Ich kann prüfen, ob eine Zahl eine Lösung einer quadratischen Gleichung ist.}
Setze die Zahl ein und kreuze an.
\begin{teile}
\teil $x = 4$ in $x^2 - 16 = 0$ \hfill \kreuz{(wA)} \quad \kreuz{(fA)}
\teil $x = -4$ in $x^2 + 16 = 0$ \hfill \kreuz{(wA)} \quad \kreuz{(fA)}
\teil $x = -3$ in $(x - 3)^2 = 36$ \hfill \kreuz{(wA)} \quad \kreuz{(fA)}
\teil $x = 3$ in $(x + 5)^2 = 4$ \hfill \kreuz{(wA)} \quad \kreuz{(fA)}
\teil $x = -2$ in $2x^2 + 1 = 9$ \hfill \kreuz{(wA)} \quad \kreuz{(fA)}
\teil $x = -1$ in $(x - 1)^2 = -4$ \hfill \kreuz{(wA)} \quad \kreuz{(fA)}
\end{teile}
\end{aufgabe}

\begin{aufgabe}{Ich kann am Graphen ablesen, wie viele Lösungen eine Gleichung hat.}
Die Grafik zeigt die Parabel zu $y = (x - 1)^2 - 2$ und drei waagerechte Geraden.

\begin{ksys}[xmin=-3,xmax=5,ymin=-4,ymax=6,ablesen]
\parabel{1}{1}{-2}{}
\gerade[4.4]{0}{2}{y=2}
\gerade[4]{0}{-2}{y=-2}
\gerade[4]{0}{-3}{y=-3}
\end{ksys}

Gib an, wie viele Lösungen die Gleichung hat.
\begin{teile}
\teil $(x - 1)^2 - 2 = 2$ \hfill Anzahl: \leerfeld
\teil $(x - 1)^2 - 2 = -2$ \hfill Anzahl: \leerfeld
\teil $(x - 1)^2 - 2 = -3$ \hfill Anzahl: \leerfeld
\teil $(x - 1)^2 - 2 = 0$ \hfill Anzahl: \leerfeld
\anweisung{Lies die Lösungen am Graphen ab.}
\teil $(x - 1)^2 - 2 = 2$ \hfill \feld{x_1} \quad \feld{x_2}
\end{teile}
\end{aufgabe}

\begin{aufgabe}{Ich kann eine Gleichung angeben, die keine Lösung hat.}
Gib eine passende Gleichung an.
\begin{teile}
\teil Form $x^2 = $ Zahl, ohne Lösung \hfill \leerfeld
\teil mit quadrierter Klammer, ohne Lösung \hfill \leerfeld
\teil mit $x^2$ und einer Zahl links, ohne Lösung \hfill \leerfeld
\end{teile}
\end{aufgabe}

\begin{aufgabe}{Ich kann Aussagen über die Lösungen einer Gleichung beurteilen.}
Kreuze an, ob die Aussage wahr oder falsch ist.
\begin{teile}
\teil Die Gleichung $(x - 5)^2 = 0$ hat genau eine Lösung. \hfill \kreuz{wahr} \quad \kreuz{falsch}
\teil $2$ und $10$ sind die Lösungen der Gleichung $(x + 6)^2 = 16$. \hfill \kreuz{wahr} \quad \kreuz{falsch}
\anweisung{Ändere die Gleichung $(x + 6)^2 = 16$ so, dass sie keine Lösung hat.}
\teil Neue rechte Seite: \leerfeld \quad (P10 2021 FOR)
\end{teile}
\end{aufgabe}

\begin{aufgabe}{Ich finde den Fehler beim Wurzelziehen.}
Finde den Fehler, benenne ihn und korrigiere die Rechnung.
\begin{teile}
\teil Ben löst $2x^2 = 162$:
\rechnung{2x^2 &= 162 &&\mid :2 \\ x^2 &= 81 \\ x &= 9}
\teil Lina löst $(x + 2)^2 = 49$:
\rechnung{x + 2 &= 7 &&\text{oder}\quad x + 2 = -7 \\ x_1 &= 9 &&\text{oder}\quad x_2 = -5}
\teil Tom löst $x^2 + 21 = 12$:
\rechnung{x^2 + 21 &= 12 &&\mid -21 \\ x^2 &= -9 \\ x &= -3}
\end{teile}
\end{aufgabe}

\begin{aufgabe}{Ich kann begründen, warum eine Gleichung zwei, eine oder keine Lösung hat.}
Begründe ohne zu rechnen.
\begin{teile}
\teil Warum hat $x^2 = 20$ zwei Lösungen?
\teil Warum hat $x^2 = -20$ keine Lösung?
\teil Tim sagt: „$(x - 4)^2 = 0$ hat keine Lösung, weil man aus $0$ keine Wurzel ziehen kann.“ Stimmt das?
\end{teile}
\end{aufgabe}

\begin{aufgabe}{Ich kann eine Länge durch Wurzelziehen berechnen.}
Stelle eine Gleichung auf, löse sie und antworte mit Einheit.
\begin{teile}
\teil Eine quadratische Terrasse hat einen Flächeninhalt von $64\,\text{m}^2$. Wie lang ist eine Seite der Terrasse? (P10 2020 FOR) \hfill Seite: \leerfeld[m]
\teil Eine zylinderförmige Regentonne ist $90\,\text{cm}$ hoch und soll $300$ Liter fassen. Für das Volumen gilt $V = \pi\cdot r^2\cdot h$ mit dem Radius $r$ und der Höhe $h$. Wie groß muss der Radius sein? Runde auf eine Stelle nach dem Komma. (P10 2022 FOR) \hfill Radius: \leerfeld[cm]
\end{teile}
\end{aufgabe}

\clearpage
% Einheit 2 – Satz vom Nullprodukt – Nr. 24–31
\setcounter{aufgabe}{23}
\einheitenkopf[e2]{Einheit 2 von 4 · Satz vom Nullprodukt}
\zweigzeile{Hier lernst du, Gleichungen der Form Produkt gleich null mit dem Satz vom Nullprodukt zu lösen · neu oder schon bekannt – je nach Buch (an der Oberschule erst Kl. 10) · P10 · baut auf: lineare Gleichungen, Ausmultiplizieren und Ausklammern, Einsetzen}

\begin{aufgabe}{Ich erkenne, ob ein Produkt gleich null gesetzt ist.}
Rechne nicht. Steht links ein Produkt und rechts null? Kreuze an.
\begin{teile}
\teil $(x - 2)\cdot(x - 9) = 0$ \hfill \janein
\teil $x\cdot(x + 4) = 12$ \hfill \janein
\teil $(x + 1)\cdot(x - 6) = 0$ \hfill \janein
\teil $(x - 3)\cdot(x + 3) = 7$ \hfill \janein
\teil $x\cdot(x + 4) - 12 = 0$ \hfill \janein
\teil $x\cdot(x - 5) = 0$ \hfill \janein
\end{teile}
\end{aufgabe}

\verfahren{Gleichungen in Produktform}

\begin{aufgabe}{Ich kann eine Gleichung in Produktform mit dem Satz vom Nullprodukt lösen.}
Löse die Gleichung.
\rechenplatz{3}
\begin{gleichungsraster}[2]
\gl{(x - 2)\cdot(x - 5) = 0} & \gl{(x - 1)\cdot(x - 4) = 0} \\
\gl{(x - 6)\cdot(x - 3) = 0} & \gl{(x - 7)\cdot(x - 2) = 0} \\
\gl{(x + 4)\cdot(x - 1) = 0} & \gl{x\cdot(x + 6) = 0} \\
\gl{(x - 3)\cdot(x - 3) = 0} & \\
\end{gleichungsraster}
\end{aufgabe}

\begin{aufgabe}{Ich kann prüfen, ob eine Zahl eine Gleichung in Produktform erfüllt.}
Setze die Zahl ein und kreuze an.
\begin{teile}
\teil $x = 3$ in $(x - 3)\cdot(x + 8) = 0$ \hfill \kreuz{(wA)} \quad \kreuz{(fA)}
\teil $x = -2$ in $(x - 2)\cdot(x + 5) = 0$ \hfill \kreuz{(wA)} \quad \kreuz{(fA)}
\teil $x = -7$ in $x\cdot(x + 9) = -14$ \hfill \kreuz{(wA)} \quad \kreuz{(fA)}
\teil $x = -1$ in $x\cdot(x - 3) = -4$ \hfill \kreuz{(wA)} \quad \kreuz{(fA)}
\teil $x = -6$ in $(x + 2)\cdot(x + 3) = 12$ \hfill \kreuz{(wA)} \quad \kreuz{(fA)}
\anweisung{Welche Zahl erfüllt die Gleichung? Kreuze an.}
\teil $x\cdot(x + 7) = -10$ \hfill \kreuz{$x = 2$} \quad \kreuz{$x = 5$} \quad \kreuz{$x = -2$} \quad \kreuz{$x = -10$} \quad (P10 2025 FOR)
\end{teile}
\end{aufgabe}

\verfahren{Gleichungen, in denen jedes Glied ein $x$ hat}

\begin{aufgabe}{Ich kann eine Gleichung, in der jedes Glied ein $x$ hat, durch Ausklammern lösen.}
Löse die Gleichung.
\rechenplatz{3}
\begin{gleichungsraster}[3]
\gl{x^2 + 4x = 0} & \gl{x^2 + 5x = 0} \\
\gl{x^2 + 2x = 0} & \gl{x^2 - 6x = 0} \\
\gl{5x^2 - 15x = 0} & \\
\end{gleichungsraster}
\end{aufgabe}

\verfahren{Zwischen Produkt und Normalform wechseln}

\begin{aufgabe}{Ich kann eine Produktgleichung mit einer Zahl rechts ausmultiplizieren und ordnen.}
Multipliziere aus und bringe alles auf die linke Seite. Löse noch nicht.
\rechenplatz{2}
\begin{gleichungsraster}[2]
\gl{x\cdot(x + 3) = 10} & \gl{x\cdot(x - 4) = -3} \\
\gl{(x + 1)\cdot(x - 2) = 4} & \\
\end{gleichungsraster}
\end{aufgabe}

\begin{aufgabe}{Ich kann eine Gleichung lösen, wenn ich ihre Produktform kenne.}
Prüfe die Produktform durch Ausmultiplizieren. Gib dann die Lösungen an.
\rechenplatz{3}
\begin{teile}
\teil $x^2 + 5x + 6 = 0$ \quad Produktform: $(x + 2)\cdot(x + 3) = 0$ \hfill \feld{x_1} \quad \feld{x_2}
\teil $x^2 - 2x - 8 = 0$ \quad Produktform: $(x - 4)\cdot(x + 2) = 0$ \hfill \feld{x_1} \quad \feld{x_2}
\teil Berechne die Nullstellen der Funktion $f$ mit $f(x) = 4x^2 - 4x - 24$. Zeige dazu, dass $4x^2 - 4x - 24 = 4\cdot(x - 3)\cdot(x + 2)$ ist. (P10 2020 FOR) \hfill \feld{x_1} \quad \feld{x_2}
\end{teile}
\end{aufgabe}

\begin{aufgabe}{Ich finde den Fehler beim Satz vom Nullprodukt.}
Finde den Fehler, benenne ihn und korrigiere die Rechnung.
\begin{teile}
\teil Jana löst $x^2 = 6x$:
\rechnung{x^2 &= 6x &&\mid :x \\ x &= 6}
\teil Paul löst $(x + 5)\cdot(x - 2) = 0$:
\rechnung{x + 5 &= 0 &&\text{oder}\quad x - 2 = 0 \\ x_1 &= 5 &&\text{oder}\quad x_2 = -2}
\teil Ole löst $x\cdot(x + 1) = 12$:
\rechnung{x &= 12 &&\text{oder}\quad x + 1 = 12 \\ x_1 &= 12 &&\text{oder}\quad x_2 = 11}
\end{teile}
\end{aufgabe}

\begin{aufgabe}{Ich kann Regeln zum Satz vom Nullprodukt begründen.}
Begründe.
\begin{teile}
\teil Warum ist ein Produkt null, sobald einer der Faktoren null ist?
\teil Warum darf man bei $x^2 = 3x$ nicht einfach durch $x$ teilen?
\end{teile}
\end{aufgabe}

\clearpage
% Einheit 3 – Normalform und p-q-Formel – Nr. 32–44
\setcounter{aufgabe}{31}
\einheitenkopf[e3]{Einheit 3 von 4 · Normalform und p-q-Formel}
\zweigzeile{Hier lernst du, eine quadratische Gleichung in die Normalform zu bringen und mit der p-q-Formel zu lösen · neu in diesem Jahr (an der Oberschule je nach Buch Kl. 9 oder 10) · P10 oft · baut auf: Wurzelziehen (Einheit 1), Ausmultiplizieren und Ordnen (Einheit 2), binomische Formeln}

\begin{aufgabe}{Ich erkenne, welche Form eine Gleichung hat und welcher Lösungsweg passt.}
Rechne nicht. Kreuze den passenden Lösungsweg an.

\sachtabelle{lcccc}{Gleichung & Wurzel ziehen & Nullprodukt & rückwärts rechnen & p-q-Formel}{a) $3x^2 = 27$ & $\square$ & $\square$ & $\square$ & $\square$ \\ b) $(x - 4)\cdot(x + 1) = 0$ & $\square$ & $\square$ & $\square$ & $\square$ \\ c) $(x + 2)^2 = 11$ & $\square$ & $\square$ & $\square$ & $\square$ \\ d) $x^2 + 5x - 14 = 0$ & $\square$ & $\square$ & $\square$ & $\square$ \\ e) $x^2 - 7x + 6 = 0$ & $\square$ & $\square$ & $\square$ & $\square$}
\end{aufgabe}

\begin{aufgabe}{Ich kann p und q mit Vorzeichen ablesen.}
Die Normalform heißt $x^2 + p\cdot x + q = 0$. Rechne nichts aus. Schreibe $p$ und $q$ mit Vorzeichen auf.
\begin{geruest}
\gz{a}{$x^2 + 9x + 4 = 0$}{\feld{p}\feld{q}}
\gz{b}{$x^2 - 5x + 2 = 0$}{\feld{p}\feld{q}}
\gz{c}{$x^2 + 2x - 8 = 0$}{\feld{p}\feld{q}}
\gz{d}{$x^2 - x - 20 = 0$}{\feld{p}\feld{q}}
\gz{e}{$x^2 - 7x = 0$}{\feld{p}\feld{q}}
\end{geruest}
\end{aufgabe}

\begin{aufgabe}{Ich erkenne, ob die Vorzahl von $x^2$ eins ist.}
Rechne nicht. Musst du vor der Formel teilen? Kreuze an und gib an, durch welche Zahl.
\begin{teile}
\teil $x^2 + 4x - 5 = 0$ \hfill \kreuz{nein} \quad \kreuz{ja, teilen durch} \leerfeld
\teil $2x^2 - 6x + 4 = 0$ \hfill \kreuz{nein} \quad \kreuz{ja, teilen durch} \leerfeld
\teil $-x^2 + 3x + 10 = 0$ \hfill \kreuz{nein} \quad \kreuz{ja, teilen durch} \leerfeld
\teil $5x^2 + 10x - 15 = 0$ \hfill \kreuz{nein} \quad \kreuz{ja, teilen durch} \leerfeld
\teil $x^2 - 3x = 0$ \hfill \kreuz{nein} \quad \kreuz{ja, teilen durch} \leerfeld
\end{teile}
\end{aufgabe}

\verfahren{Lösen mit der p-q-Formel}

\begin{aufgabe}{Ich kann eine Gleichung in Normalform mit der p-q-Formel lösen.}
Löse mit der p-q-Formel.
\rechenplatz{4}
\begin{gleichungsraster}[4]
\gl{x^2 + 6x + 8 = 0} & \gl{x^2 + 10x + 21 = 0} \\
\gl{x^2 + 12x + 20 = 0} & \gl{x^2 + 6x + 5 = 0} \\
\gl{x^2 - 2x - 8 = 0} & \\
\end{gleichungsraster}
\end{aufgabe}

\begin{aufgabe}{Ich kann eine Gleichung erst ordnen oder normieren und dann mit der Formel lösen.}
Bringe die Gleichung in die Normalform und löse sie.
\rechenplatz{2}
\begin{gleichungsraster}[4]
\gl{x^2 + 5 = 6x} & \gl{x^2 = 2x + 3} \\
\gl{2x^2 + 4x - 16 = 0} & \gl{-x^2 + 6x - 8 = 0} \\
\gl{2x^2 + 2x - 12 = 0} & \\
\end{gleichungsraster}
\end{aufgabe}

\begin{aufgabe}{Ich kann mit der Formel erkennen, ob eine Gleichung zwei, eine oder keine Lösung hat.}
Löse mit der p-q-Formel.
\rechenplatz{2}
\begin{gleichungsraster}[3]
\gl{x^2 - 6x + 9 = 0} & \gl{x^2 - 2x + 5 = 0} \\
\anweisung{Berechne nur den Wert unter der Wurzel. Er heißt Diskriminante. Gib an, wie viele Lösungen die Gleichung hat.}
\gl[2]{x^2 - 3x + 2 = 0} & \gl[2]{x^2 + 5x + 7 = 0} \\
\gl[2]{x^2 - 7x + 12{,}25 = 0} & \\
\end{gleichungsraster}
\end{aufgabe}

\begin{aufgabe}{Ich kann Lösungen mit Wurzel und als Näherungswert angeben.}
Löse mit der p-q-Formel. Gib die Lösungen mit Wurzel und auf zwei Stellen nach dem Komma gerundet an.
\rechenplatz{2}
\begin{gleichungsraster}[4]
\gl{x^2 - 4x + 1 = 0} & \gl{x^2 + x - 1 = 0} \\
\gl[5]{3x^2 - 6x - 3 = 0} & \\
\end{gleichungsraster}
\end{aufgabe}

\begin{aufgabe}{Ich kann mit einer Lösung die Probe machen.}
Löse die Gleichung und mache mit einer Lösung die Probe.
\rechenplatz{2}
\begin{gleichungsraster}[5]
\gl{x^2 - 2x - 24 = 0} & \gl{3x^2 - 12x - 15 = 0} \\
\end{gleichungsraster}
\end{aufgabe}

\verfahren{Gleichungen mit Klammern}

\begin{aufgabe}{Ich kann eine Gleichung mit Klammern erst auflösen, dann ordnen und lösen.}
Löse die Gleichung.
\rechenplatz{3}
\begin{gleichungsraster}[5]
\gl{x\cdot(x + 3) = 18} & \gl{(x + 1)^2 = 3x + 7} \\
\gl{(x - 2)\cdot(x + 4) = 3x + 4} & \\
\end{gleichungsraster}
\end{aufgabe}

\verfahren{Den Lösungsweg wählen}

\begin{aufgabe}{Ich kann eine Gleichung auf dem passenden Weg lösen.}
Löse die Gleichung auf einem geschickten Weg.
\begin{gleichungsraster}[3]
\gl{4x^2 = 100} & \gl{x^2 + 7x = 0} \\
\gl{(x - 4)^2 = 1} & \gl[4]{x^2 - 3x - 4 = 0} \\
\end{gleichungsraster}
\begin{teile}
\teil Für die Parabel mit $f(x) = -x^2 + 4x + 9$ sollen alle Stellen $x$ berechnet werden, an denen $f(x) = -3$ gilt. Berechne sie. (P10 2023 FOR) \hfill \feld{x_1} \quad \feld{x_2}
\end{teile}
\end{aufgabe}

\begin{aufgabe}{Ich finde den Fehler bei der p-q-Formel.}
Finde den Fehler, benenne ihn und korrigiere die Rechnung.
\begin{teile}
\teil Emma löst $3x^2 + 6x - 9 = 0$:
\rechnung{p &= 6, \quad q = -9 \\ x_{1,2} &= -3 \pm \sqrt{9 + 9} \\ &= -3 \pm \sqrt{18}}
\teil Felix löst $x^2 - 8x + 12 = 0$:
\rechnung{x_{1,2} &= -4 \pm \sqrt{16 - 12} \\ &= -4 \pm 2 \\ x_1 &= -2, \quad x_2 = -6}
\teil Nora löst $x^2 - 10x + 16 = 0$:
\rechnung{x_{1,2} &= 5 \pm \sqrt{25 + 16} \\ &= 5 \pm \sqrt{41}}
\end{teile}
\end{aufgabe}

\begin{aufgabe}{Ich kann Regeln der p-q-Formel begründen.}
Begründe.
\begin{teile}
\teil Warum hat eine Gleichung keine Lösung, wenn unter der Wurzel eine negative Zahl steht?
\teil Warum muss man $2x^2 - 10x + 12 = 0$ erst durch 2 teilen, bevor man die p-q-Formel benutzt?
\end{teile}
\end{aufgabe}

\begin{aufgabe}{Ich kann die Schnittpunkte einer Geraden mit einer Parabel berechnen.}
Die Grafik zeigt die Parabel zu $f(x) = x^2 - 2$ und die Gerade zu $g(x) = -2x + 1$.

\begin{ksys}[xmin=-4,xmax=3,ymin=-3,ymax=8,ablesen]
\parabel{1}{0}{-2}{f}
\gerade{-2}{1}{g}
\end{ksys}

\begin{teile}
\teil Setze $f(x) = g(x)$ und berechne die Schnittpunkte. \hfill \feld{S_1} \quad \feld{S_2}
\teil Lies die Schnittpunkte am Graphen ab und vergleiche mit a).
\teil Gegeben sind die Parabel mit $f(x) = (x - 1)^2 - 4$ und die Gerade mit $g(x) = 3x + 3$. Berechne die Koordinaten der Schnittpunkte. (P10 2022 FOR) \hfill \feld{S_1} \quad \feld{S_2}
\end{teile}
\end{aufgabe}

\clearpage
% Einheit 4 – Sachaufgaben – Nr. 45–49
\setcounter{aufgabe}{44}
\einheitenkopf[e4]{Einheit 4 von 4 · Sachaufgaben}
\zweigzeile{Hier lernst du, zu Zahlenrätseln und Rechtecken eine quadratische Gleichung aufzustellen, sie zu lösen und die passende Lösung auszuwählen · neu in diesem Jahr (an der Oberschule je nach Buch Kl. 9 oder 10) · keine P10-Aufgabe · baut auf: Einheit 1 bis 3, Gleichungen aufstellen}

\begin{aufgabe}{Ich erkenne, welche Lösung zur Frage passt.}
Rechne nicht. Eine Gleichung zur Frage hat die Lösungen $x_1$ und $x_2$. Kreuze an, welche Lösung als Antwort passt.
\begin{teile}
\teil Seitenlänge eines quadratischen Beets in m: $x_1 = 6$, $x_2 = -6$ \hfill \kreuz{$x_1$} \quad \kreuz{$x_2$} \quad \kreuz{beide}
\teil Breite eines Rechtecks in cm: $x_1 = -12$, $x_2 = 4$ \hfill \kreuz{$x_1$} \quad \kreuz{$x_2$} \quad \kreuz{beide}
\teil Anzahl der Stuhlreihen in einem Saal: $x_1 = 9$, $x_2 = -14$ \hfill \kreuz{$x_1$} \quad \kreuz{$x_2$} \quad \kreuz{beide}
\teil Zahlen, deren Quadrat 144 ist: $x_1 = 12$, $x_2 = -12$ \hfill \kreuz{$x_1$} \quad \kreuz{$x_2$} \quad \kreuz{beide}
\teil Alter eines Kindes in Jahren: $x_1 = -3$, $x_2 = 11$ \hfill \kreuz{$x_1$} \quad \kreuz{$x_2$} \quad \kreuz{beide}
\end{teile}
\end{aufgabe}

\verfahren{Zahlenrätsel}

\begin{aufgabe}{Ich kann ein Zahlenrätsel mit einer quadratischen Gleichung lösen.}
Löse das Rätsel. Gib alle passenden Zahlen an.
\begin{teile}
\teil Das Quadrat einer Zahl ist 100. \quad Gleichung: $x^2 = 100$ \hfill Zahlen: \leerfeld
\teil Das Quadrat einer Zahl ist 121. \quad Gleichung: $x^2 = 121$ \hfill Zahlen: \leerfeld
\teil Das Doppelte des Quadrats einer Zahl ist 72. \quad Gleichung: $2x^2 = 72$ \hfill Zahlen: \leerfeld
\teil Vermindert man das Quadrat einer Zahl um 10, erhält man 71. \\ Gleichung: $x^2 - 10 = 71$ \hfill Zahlen: \leerfeld
\anweisung{Stelle die Gleichung selbst auf.}
\teil Addiert man 20 zum Quadrat einer Zahl, erhält man 84. \\ Gleichung: \leerfeld \hfill Zahlen: \leerfeld
\teil Das Produkt einer Zahl und ihres Nachfolgers ist 56. \\ Gleichung: \leerfeld \hfill Zahlenpaare: \leerfeld
\end{teile}
\end{aufgabe}

\verfahren{Rechteckaufgaben}

\begin{aufgabe}{Ich kann eine Rechteckaufgabe mit einer quadratischen Gleichung lösen.}
Ein rechteckiges Beet ist 3\,m länger als breit und hat einen Flächeninhalt von $40\,\text{m}^2$.

\rechteck[punkte=,seiten={x + 3,x,,}]{4.8}{3}

\begin{teile}
\teil Stelle eine Gleichung für den Flächeninhalt auf. \hfill \leerfeld
\teil Löse die Gleichung. \hfill \feld{x_1} \quad \feld{x_2}
\teil Begründe, welche Lösung zur Aufgabe passt.
\teil Gib Länge und Breite des Beets in einem Antwortsatz mit Einheit an.
\teil Prüfe am Beet, ob der Flächeninhalt stimmt.
\end{teile}
Um ein anderes rechteckiges Beet, das 4\,m länger als breit ist, wird ringsum ein 1\,m breiter Weg angelegt. Beet und Weg bedecken zusammen $96\,\text{m}^2$.

\rechteckrand{2}{3.6}{0.6}{x}{x + 4}{Weg, 1\,m breit}

\begin{teile}
\teil Berechne die Länge und die Breite des Beets. \hfill \feld[m]{\text{Breite}} \quad \feld[m]{\text{Länge}}
\end{teile}
\end{aufgabe}

\begin{aufgabe}{Ich finde den Fehler bei einer Sachaufgabe.}
Finde den Fehler, benenne ihn und korrigiere.
\begin{teile}
\teil Ein Rechteck ist 2\,cm länger als breit und hat einen Flächeninhalt von $48\,\text{cm}^2$. Sven rechnet:
\rechnung{x\cdot(x + 2) &= 48 \\ x^2 + 2x - 48 &= 0 \\ x_1 = 6, \quad x_2 &= -8}
Svens Antwort: Das Rechteck ist $-8$\,cm breit und $-6$\,cm lang.
\teil Ein Rechteck ist 5\,m länger als breit und hat einen Flächeninhalt von $24\,\text{m}^2$. Lisa rechnet:
\rechnung{2x + 2\cdot(x + 5) &= 24 \\ 4x + 10 &= 24 \\ x &= 3{,}5}
\end{teile}
\end{aufgabe}

\begin{aufgabe}{Ich kann begründen, wann eine Sachaufgabe eine oder zwei Antworten hat.}
Begründe.
\begin{teile}
\teil Warum passt bei einer Rechteckaufgabe meist nur eine Lösung der Gleichung?
\teil Warum hat das Rätsel „Das Quadrat einer Zahl ist 25. Wie heißt die Zahl?“ zwei Antworten?
\end{teile}
\end{aufgabe}

% Abhakseite „Das kann ich" – erzeugt von abhaken_bau.py aus den Quelltexten
\begin{abhakseite}
\abhakgruppe{Kennst du schon}
\abhak{1}{Ich kann Zahlen quadrieren, auch negative Zahlen und Dezimalzahlen.}
\abhak{2}{Ich kann mit Klammern und negativen Zahlen rechnen.}
\abhak{3}{Ich kann durch Einsetzen prüfen, ob eine Zahl eine Gleichung erfüllt.}
\abhak{4}{Ich finde den Fehler beim Einsetzen einer negativen Zahl.}
\abhak{5}{Ich kann eine negative Zahl in einen Term mit Quadrat einsetzen.}
\abhak{6}{Ich kann eine lineare Gleichung lösen.}
\abhak{7}{Ich kann eine Quadratwurzel ziehen.}
\abhak{8}{Ich kann den Scheitelpunkt aus der Scheitelpunktform ablesen.}
\abhak{9}{Ich kann Klammern ausmultiplizieren, auch mit binomischen Formeln.}
\abhak{10}{Ich kann $x$ ausklammern.}
\abhak{11}{Ich kann Terme ordnen und zusammenfassen.}
\abhakgruppe{Einheit 1 · Wurzelziehen und Lösbarkeit}
\abhak{12}{Ich erkenne, ob eine Gleichung quadratisch ist.}
\abhak{13}{Ich erkenne an der Zahl rechts, wie viele Lösungen eine Gleichung wie $x^2 = 5$ hat.}
\abhakgruppe{– Gleichungen der Form $x^2 = $ Zahl}
\abhak{14}{Ich kann eine Gleichung der Form $x^2 = $ Zahl durch Wurzelziehen lösen.}
\abhak{15}{Ich kann erst $x^2$ freistellen und dann die Wurzel ziehen.}
\abhakgruppe{– Gleichungen mit quadrierter Klammer}
\abhak{16}{Ich kann eine Gleichung mit quadrierter Klammer durch Rückwärtsrechnen lösen.}
\abhakgruppe{– Lösungen prüfen und zählen}
\abhak{17}{Ich kann prüfen, ob eine Zahl eine Lösung einer quadratischen Gleichung ist.}
\abhak{18}{Ich kann am Graphen ablesen, wie viele Lösungen eine Gleichung hat.}
\abhak{19}{Ich kann eine Gleichung angeben, die keine Lösung hat.}
\abhak{20}{Ich kann Aussagen über die Lösungen einer Gleichung beurteilen.}
\abhak{21}{Ich finde den Fehler beim Wurzelziehen.}
\abhak{22}{Ich kann begründen, warum eine Gleichung zwei, eine oder keine Lösung hat.}
\abhak{23}{Ich kann eine Länge durch Wurzelziehen berechnen.}
\abhakgruppe{Einheit 2 · Satz vom Nullprodukt}
\abhak{24}{Ich erkenne, ob ein Produkt gleich null gesetzt ist.}
\abhakgruppe{– Gleichungen in Produktform}
\abhak{25}{Ich kann eine Gleichung in Produktform mit dem Satz vom Nullprodukt lösen.}
\abhak{26}{Ich kann prüfen, ob eine Zahl eine Gleichung in Produktform erfüllt.}
\abhakgruppe{– Gleichungen, in denen jedes Glied ein $x$ hat}
\abhak{27}{Ich kann eine Gleichung, in der jedes Glied ein $x$ hat, durch Ausklammern lösen.}
\abhakgruppe{– Zwischen Produkt und Normalform wechseln}
\abhak{28}{Ich kann eine Produktgleichung mit einer Zahl rechts ausmultiplizieren und ordnen.}
\abhak{29}{Ich kann eine Gleichung lösen, wenn ich ihre Produktform kenne.}
\abhak{30}{Ich finde den Fehler beim Satz vom Nullprodukt.}
\abhak{31}{Ich kann Regeln zum Satz vom Nullprodukt begründen.}
\abhakgruppe{Einheit 3 · Normalform und p-q-Formel}
\abhak{32}{Ich erkenne, welche Form eine Gleichung hat und welcher Lösungsweg passt.}
\abhak{33}{Ich kann p und q mit Vorzeichen ablesen.}
\abhak{34}{Ich erkenne, ob die Vorzahl von $x^2$ eins ist.}
\abhakgruppe{– Lösen mit der p-q-Formel}
\abhak{35}{Ich kann eine Gleichung in Normalform mit der p-q-Formel lösen.}
\abhak{36}{Ich kann eine Gleichung erst ordnen oder normieren und dann mit der Formel lösen.}
\abhak{37}{Ich kann mit der Formel erkennen, ob eine Gleichung zwei, eine oder keine Lösung hat.}
\abhak{38}{Ich kann Lösungen mit Wurzel und als Näherungswert angeben.}
\abhak{39}{Ich kann mit einer Lösung die Probe machen.}
\abhakgruppe{– Gleichungen mit Klammern}
\abhak{40}{Ich kann eine Gleichung mit Klammern erst auflösen, dann ordnen und lösen.}
\abhakgruppe{– Den Lösungsweg wählen}
\abhak{41}{Ich kann eine Gleichung auf dem passenden Weg lösen.}
\abhak{42}{Ich finde den Fehler bei der p-q-Formel.}
\abhak{43}{Ich kann Regeln der p-q-Formel begründen.}
\abhak{44}{Ich kann die Schnittpunkte einer Geraden mit einer Parabel berechnen.}
\abhakgruppe{Einheit 4 · Sachaufgaben}
\abhak{45}{Ich erkenne, welche Lösung zur Frage passt.}
\abhakgruppe{– Zahlenrätsel}
\abhak{46}{Ich kann ein Zahlenrätsel mit einer quadratischen Gleichung lösen.}
\abhakgruppe{– Rechteckaufgaben}
\abhak{47}{Ich kann eine Rechteckaufgabe mit einer quadratischen Gleichung lösen.}
\abhak{48}{Ich finde den Fehler bei einer Sachaufgabe.}
\abhak{49}{Ich kann begründen, wann eine Sachaufgabe eine oder zwei Antworten hat.}
\end{abhakseite}

\end{document}
